\documentclass[11pt,a4paper]{article}

% ---------- Encodage et langue ----------
\usepackage[utf8]{inputenc}
\usepackage[T1]{fontenc}
\usepackage[french]{babel}

% ---------- Mise en page ----------
\usepackage[a4paper,margin=2cm]{geometry}
\usepackage{lmodern}
\usepackage{microtype}
\usepackage{setspace}
\setstretch{1.08}

% ---------- Maths ----------
\usepackage{amsmath,amssymb,amsthm,mathrsfs}
\usepackage{bm}

% ---------- Couleurs et encadrés ----------
\usepackage{xcolor}
\usepackage[most]{tcolorbox}

% ---------- Listes ----------
\usepackage{enumitem}

% ---------- Liens ----------
\usepackage{hyperref}
\hypersetup{
    colorlinks=true,
    linkcolor=blue!50!black,
    urlcolor=blue!50!black
}

% ---------- En-têtes / pieds ----------
\usepackage{fancyhdr}
\pagestyle{fancy}
\fancyhf{}
\fancyhead[L]{Banque d'exercices --- Terminale spécialité mathématiques}
\fancyhead[R]{Raisonnements poussés}
\fancyfoot[C]{\thepage}

% ---------- Commandes ----------
\newcommand{\R}{\mathbb{R}}
\newcommand{\C}{\mathbb{C}}
\newcommand{\N}{\mathbb{N}}

% ---------- Style des sections ----------
\definecolor{bleuclair}{RGB}{225,238,255}
\definecolor{bleufonce}{RGB}{25,70,140}
\definecolor{grisclair}{RGB}{245,245,245}

\newtcolorbox{chapitrebox}[1]{
    colback=bleuclair,
    colframe=bleufonce,
    boxrule=0.8pt,
    arc=2mm,
    left=2mm,right=2mm,top=1mm,bottom=1mm,
    title=\bfseries #1
}

\newtcolorbox{objectifbox}{
    colback=grisclair,
    colframe=black!50,
    boxrule=0.5pt,
    arc=1.5mm,
    left=2mm,right=2mm,top=1mm,bottom=1mm
}

% ---------- Document ----------
\begin{document}

\begin{center}
    {\LARGE \textbf{Banque d'exercices de Terminale}}\\[0.3em]
    {\large \textbf{avec raisonnements légèrement plus poussés}}\\[0.5em]
    {\normalsize Niveau Terminale.
\end{center}

\vspace{0.5em}

\begin{objectifbox}
\textbf{Objectif.} Cette feuille regroupe des exercices classés par chapitre.  
Les outils utilisés restent ceux du \textbf{programme de Terminale}, mais les raisonnements demandés sont plus fins, plus structurés et plus rigoureux.
\end{objectifbox}

\vspace{1em}

% =========================================================
\begin{chapitrebox}{Chapitre 1 --- Fonctions}
\end{chapitrebox}

\begin{enumerate}[label=\textbf{Exercice \arabic*.}, leftmargin=*, itemsep=1em]

\item Soit \(f(x)=\dfrac{x^2-1}{x-1}\) définie sur \(\R\setminus\{1\}\).

\begin{enumerate}[label=\alph*)]
    \item Simplifier \(f(x)\).
    \item Étudier les variations de \(f\).
    \item Peut-on prolonger \(f\) par continuité en \(x=1\) ?
    \item Interpréter graphiquement le résultat.
\end{enumerate}

\item Comparer \(e^\pi\) et \(\pi^e\).

\item Soit \(f(x)=\dfrac{x}{x^2+1}\).

\begin{enumerate}[label=\alph*)]
    \item Étudier les variations de \(f\) sur \(\R\).
    \item Montrer que, pour tout réel \(x\),
    \[
    -\frac12 \leq f(x)\leq \frac12.
    \]
    \item Retrouver ce résultat sans utiliser la dérivée.
\end{enumerate}

\item Soit \(f(x)=\sqrt{x+2}-\sqrt{x}\) définie sur \([0,+\infty[\).

\begin{enumerate}[label=\alph*)]
    \item Étudier les variations de \(f\).
    \item Déterminer sa limite en \(+\infty\).
    \item Montrer que \(f(x)>0\) pour tout \(x\geq 0\).
    \item Déterminer le plus grand réel \(M\) tel que \(f(x)\leq M\) pour tout \(x\geq 0\).
\end{enumerate}

\item Déterminer le nombre de solutions de l'équation
\[
x+\sqrt{x}=2.
\]

\begin{enumerate}[label=\alph*)]
    \item Résoudre algébriquement.
    \item Retrouver le résultat par une étude de fonction.
    \item Expliquer l'intérêt de croiser les deux méthodes.
\end{enumerate}

\end{enumerate}

% =========================================================
\begin{chapitrebox}{Chapitre 2 --- Suites}
\end{chapitrebox}

\begin{enumerate}[label=\textbf{Exercice \arabic*.}, start=1, leftmargin=*, itemsep=1em]

\item Soit \((u_n)\) définie par
\[
u_0=1,\qquad u_{n+1}=\frac{2u_n+3}{u_n+4}.
\]

\begin{enumerate}[label=\alph*)]
    \item Montrer que la suite est bien définie.
    \item Conjecturer sa limite.
    \item Montrer que si \(u_n\to \ell\), alors \(\ell\) vérifie une équation.
    \item Montrer que \((u_n)\) est convergente.
\end{enumerate}

\item Soit \((u_n)\) définie par
\[
u_0\in ]0,1[,\qquad u_{n+1}=u_n(2-u_n).
\]

\begin{enumerate}[label=\alph*)]
    \item Montrer que \(0<u_n<1\) pour tout \(n\).
    \item Étudier le sens de variation de \((u_n)\).
    \item Déterminer la limite de \((u_n)\).
\end{enumerate}

\item Soit
\[
u_n=\sum_{k=1}^n \frac{1}{k(k+1)}.
\]

\begin{enumerate}[label=\alph*)]
    \item Calculer les premiers termes.
    \item Simplifier cette somme.
    \item Déterminer la limite de \((u_n)\).
\end{enumerate}

\item Soit \((u_n)\) définie par
\[
u_0=1,\qquad u_{n+1}=\sqrt{2+u_n}.
\]

\begin{enumerate}[label=\alph*)]
    \item Montrer que \(u_n\leq 2\) pour tout \(n\).
    \item Étudier la monotonie de la suite.
    \item Montrer qu'elle converge et déterminer sa limite.
\end{enumerate}

\item Soit
\[
u_n=\frac{n}{2^n}.
\]

\begin{enumerate}[label=\alph*)]
    \item Étudier le sens de variation de \((u_n)\).
    \item Déterminer sa limite.
    \item Montrer qu'il existe un entier \(N\) tel que, pour tout \(n\geq N\),
    \[
    u_n\leq 10^{-3}.
    \]
\end{enumerate}

\end{enumerate}

% =========================================================
\begin{chapitrebox}{Chapitre 3 --- Dérivation et convexité}
\end{chapitrebox}

\begin{enumerate}[label=\textbf{Exercice \arabic*.}, leftmargin=*, itemsep=1em]

\item Soit \(f(x)=x^3-3x\).

\begin{enumerate}[label=\alph*)]
    \item Étudier les variations de \(f\).
    \item Étudier la convexité de \(f\).
    \item Déterminer les tangentes horizontales.
    \item Déterminer les points d'inflexion.
\end{enumerate}

\item Soit \(f(x)=\ln x\) sur \(]0,+\infty[\).

\begin{enumerate}[label=\alph*)]
    \item Montrer que la courbe de \(f\) est en dessous de chacune de ses tangentes.
    \item En déduire que, pour tout \(x>0\),
    \[
    \ln x\leq x-1.
    \]
\end{enumerate}

\item Soit \(f(x)=e^x\).

\begin{enumerate}[label=\alph*)]
    \item Écrire l'équation de la tangente au point d'abscisse \(0\).
    \item Comparer \(e^x\) et \(x+1\).
    \item Étudier le signe de \(e^x-x-1\).
\end{enumerate}

\item Soit \(f(x)=\dfrac{1}{x}\) sur \(]0,+\infty[\).

\begin{enumerate}[label=\alph*)]
    \item Étudier la convexité de \(f\).
    \item Montrer que, pour tous réels \(a,b>0\),
    \[
    \frac{1}{\frac{a+b}{2}} \leq \frac12\left(\frac1a+\frac1b\right).
    \]
\end{enumerate}

\item Soit \(f(x)=x\ln x\) sur \(]0,+\infty[\).

\begin{enumerate}[label=\alph*)]
    \item Étudier les variations de \(f\).
    \item Étudier sa convexité.
    \item Montrer que l'équation \(x\ln x=1\) admet une unique solution sur \(]1,+\infty[\).
\end{enumerate}

\end{enumerate}

% =========================================================
\begin{chapitrebox}{Chapitre 4 --- Logarithme et exponentielle}
\end{chapitrebox}

\begin{enumerate}[label=\textbf{Exercice \arabic*.}, leftmargin=*, itemsep=1em]

\item Résoudre dans \(\R\) :
\[
e^{2x}-3e^x+2=0.
\]

\item Résoudre :
\[
\ln(x+1)+\ln(x-2)=\ln 4.
\]

\begin{enumerate}[label=\alph*)]
    \item Déterminer l'ensemble de définition.
    \item Résoudre l'équation.
    \item Vérifier les solutions obtenues.
\end{enumerate}

\item Étudier, selon le réel \(m\), le nombre de solutions de l'équation
\[
\ln x=mx.
\]

\item Étudier les variations de la fonction
\[
f(x)=\frac{\ln x}{x}
\qquad \text{sur } ]0,+\infty[.
\]
En déduire une nouvelle démonstration pour comparer \(e^\pi\) et \(\pi^e\).

\item Montrer que, pour tout \(x>0\),
\[
\ln x\leq x-1,
\]
puis discuter le cas d'égalité.

\end{enumerate}

% =========================================================
\begin{chapitrebox}{Chapitre 5 --- Intégration et primitives}
\end{chapitrebox}

\begin{enumerate}[label=\textbf{Exercice \arabic*.}, leftmargin=*, itemsep=1em]

\item Calculer
\[
\int_0^1 (2x+3)\,dx,
\]
puis retrouver le résultat par une interprétation géométrique.

\item Calculer
\[
\int_0^1 xe^x\,dx.
\]

\item Soit
\[
f(x)=\frac{1}{x^2+1}.
\]

\begin{enumerate}[label=\alph*)]
    \item Montrer que, pour tout \(x\in[0,1]\),
    \[
    \frac12\leq f(x)\leq 1.
    \]
    \item En déduire un encadrement de
    \[
    \int_0^1 \frac{1}{x^2+1}\,dx.
    \]
\end{enumerate}

\item On considère
\[
A(x)=\int_0^x (t^2-2t+3)\,dt.
\]

\begin{enumerate}[label=\alph*)]
    \item Déterminer \(A(x)\).
    \item Étudier ses variations.
    \item Interpréter \(A(x)\) comme une aire algébrique.
\end{enumerate}

\item Pour tout entier \(n\in\N\), on pose
\[
I_n=\int_0^1 x^n\,dx.
\]

\begin{enumerate}[label=\alph*)]
    \item Calculer \(I_n\).
    \item Étudier la suite \((I_n)\).
    \item Déterminer sa limite.
\end{enumerate}

\end{enumerate}

% =========================================================
\begin{chapitrebox}{Chapitre 6 --- Équations différentielles}
\end{chapitrebox}

\begin{enumerate}[label=\textbf{Exercice \arabic*.}, leftmargin=*, itemsep=1em]

\item Résoudre sur \(\R\) :
\[
y'+2y=0.
\]

\item Déterminer la solution de
\[
y'-3y=0
\]
vérifiant \(y(0)=5\).

\item Résoudre
\[
y'+y=2.
\]

\item On considère l'équation
\[
y'=ky.
\]
Déterminer, selon le réel \(k\), les solutions croissantes.

\item Une population bactérienne est modélisée par
\[
P'(t)=0{,}3P(t),\qquad P(0)=500.
\]

\begin{enumerate}[label=\alph*)]
    \item Déterminer \(P(t)\).
    \item À quel instant la population dépasse-t-elle 2000 ?
    \item Interpréter le modèle.
\end{enumerate}

\end{enumerate}

% =========================================================
\begin{chapitrebox}{Chapitre 7 --- Probabilités}
\end{chapitrebox}

\begin{enumerate}[label=\textbf{Exercice \arabic*.}, leftmargin=*, itemsep=1em]

\item On lance deux dés équilibrés.

\begin{enumerate}[label=\alph*)]
    \item Déterminer la probabilité d'obtenir une somme égale à 7.
    \item Déterminer la probabilité d'obtenir une somme supérieure ou égale à 10.
    \item Les événements précédents sont-ils incompatibles ?
\end{enumerate}

\item Dans une urne, on place 5 boules rouges, 3 bleues et 2 vertes.  
On tire une boule au hasard.

\begin{enumerate}[label=\alph*)]
    \item Calculer la probabilité d'obtenir une boule rouge.
    \item Calculer la probabilité de ne pas obtenir une boule verte.
    \item Calculer la probabilité d'obtenir une boule rouge ou bleue.
\end{enumerate}

\item On sait que
\[
P(A)=0{,}6,\qquad P(B)=0{,}5,\qquad P(A\cap B)=0{,}3.
\]

\begin{enumerate}[label=\alph*)]
    \item Calculer \(P(A\cup B)\).
    \item Les événements \(A\) et \(B\) sont-ils indépendants ?
    \item Calculer \(P_A(B)\).
\end{enumerate}

\item Dans une population, \(2\%\) des individus ont une maladie.  
Un test est positif dans \(95\%\) des cas chez les malades, et dans \(4\%\) des cas chez les non-malades.

\begin{enumerate}[label=\alph*)]
    \item Construire un arbre pondéré.
    \item Calculer la probabilité qu'un individu soit testé positif.
    \item Calculer la probabilité qu'un individu soit malade sachant que le test est positif.
\end{enumerate}

\item On considère deux événements \(A\) et \(B\) tels que \(P(B)\neq 0\).

\begin{enumerate}[label=\alph*)]
    \item Montrer que
    \[
    P(A\cap B)=P(B)\,P_B(A).
    \]
    \item Retrouver la formule des probabilités totales dans un cas simple.
    \item Donner un exemple numérique.
\end{enumerate}

\end{enumerate}

% =========================================================
\begin{chapitrebox}{Chapitre 8 --- Variables aléatoires et loi binomiale}
\end{chapitrebox}

\begin{enumerate}[label=\textbf{Exercice \arabic*.}, leftmargin=*, itemsep=1em]

\item On lance 5 fois une pièce équilibrée.  
Soit \(X\) le nombre de piles obtenus.

\begin{enumerate}[label=\alph*)]
    \item Déterminer la loi de \(X\).
    \item Calculer \(P(X=3)\).
    \item Calculer \(P(X\geq 1)\).
\end{enumerate}

\item Une machine produit des pièces avec \(3\%\) de défauts.  
On prélève 20 pièces au hasard. Soit \(X\) le nombre de pièces défectueuses.

\begin{enumerate}[label=\alph*)]
    \item Justifier que \(X\) suit une loi binomiale.
    \item Déterminer ses paramètres.
    \item Calculer \(P(X=0)\).
    \item Calculer \(P(X\leq 2)\).
\end{enumerate}

\item On considère \(X\sim\mathcal{B}(n,p)\).

\begin{enumerate}[label=\alph*)]
    \item Donner l'espérance \(E(X)\).
    \item Dans le cas \(n=10\) et \(p=0{,}2\), calculer \(E(X)\).
    \item Interpréter le résultat.
\end{enumerate}

\item On joue à un jeu où la probabilité de gagner une partie est \(0{,}4\).  
On joue 8 parties. Soit \(X\) le nombre de victoires.

\begin{enumerate}[label=\alph*)]
    \item Calculer \(P(X=4)\).
    \item Calculer \(P(X\geq 4)\).
    \item Déterminer le nombre moyen de victoires.
\end{enumerate}

\item Une variable aléatoire \(X\) suit une loi binomiale \(\mathcal{B}(6,p)\).  
On sait que
\[
P(X=0)=\frac1{64}.
\]

\begin{enumerate}[label=\alph*)]
    \item Déterminer \(p\).
    \item Calculer \(P(X=2)\).
    \item Calculer l'espérance de \(X\).
\end{enumerate}

\end{enumerate}

% =========================================================
\begin{chapitrebox}{Chapitre 9 --- Géométrie dans l'espace}
\end{chapitrebox}

\begin{enumerate}[label=\textbf{Exercice \arabic*.}, leftmargin=*, itemsep=1em]

\item Dans un repère de l'espace, on considère
\[
A(1,0,2),\quad B(2,1,3),\quad C(0,1,1).
\]

\begin{enumerate}[label=\alph*)]
    \item Déterminer \(\overrightarrow{AB}\) et \(\overrightarrow{AC}\).
    \item Montrer que les points \(A\), \(B\) et \(C\) ne sont pas alignés.
\end{enumerate}

\item Déterminer une représentation paramétrique de la droite passant par \(A(1,2,3)\) et de vecteur directeur \((2,-1,4)\).

\item On considère le plan \((P)\) d'équation
\[
2x-y+z-3=0.
\]

\begin{enumerate}[label=\alph*)]
    \item Vérifier que \(A(1,0,1)\) appartient à \((P)\).
    \item Donner un vecteur normal au plan.
    \item Déterminer si la droite de vecteur directeur \((2,-1,1)\) est parallèle au plan.
\end{enumerate}

\item On considère les droites
\[
d:\begin{cases}
x=1+t\\
y=2-t\\
z=3+2t
\end{cases}
\qquad
d':\begin{cases}
x=2+s\\
y=1+2s\\
z=5+s
\end{cases}
\]

\begin{enumerate}[label=\alph*)]
    \item Déterminer si elles sont parallèles.
    \item Déterminer si elles sont sécantes.
    \item Conclure sur leur position relative.
\end{enumerate}

\item Dans un cube, montrer qu'une droite est perpendiculaire à un plan à l'aide d'un raisonnement vectoriel.

\end{enumerate}

% =========================================================
\begin{chapitrebox}{Chapitre 10 --- Nombres complexes}
\end{chapitrebox}

\begin{enumerate}[label=\textbf{Exercice \arabic*.}, leftmargin=*, itemsep=1em]

\item Résoudre dans \(\C\) :
\[
z^2+4=0.
\]

\item Mettre sous forme algébrique :
\[
(2+3i)(1-4i).
\]

\item Résoudre dans \(\C\) :
\[
z^2-2z+5=0.
\]

\item Soit \(z=1+i\sqrt{3}\).

\begin{enumerate}[label=\alph*)]
    \item Déterminer le module de \(z\).
    \item Déterminer un argument de \(z\).
    \item Écrire \(z\) sous forme trigonométrique.
\end{enumerate}

\item On considère les points \(A\), \(B\), \(C\) d'affixes respectives \(1\), \(i\), \(1+i\).

\begin{enumerate}[label=\alph*)]
    \item Placer ces points dans le plan complexe.
    \item Montrer que le triangle \(ABC\) est rectangle.
    \item Montrer qu'il est isocèle.
\end{enumerate}

\end{enumerate}

% =========================================================
\begin{chapitrebox}{Chapitre 11 --- Algorithmique et raisonnement}
\end{chapitrebox}

\begin{enumerate}[label=\textbf{Exercice \arabic*.}, leftmargin=*, itemsep=1em]

\item On considère l'algorithme suivant :

\begin{center}
\begin{tabular}{l}
\(u\leftarrow 1\)\\
Pour \(k\) allant de \(1\) à \(n\)\\
\hspace{1em} \(u\leftarrow 2u+1\)
\end{tabular}
\end{center}

\begin{enumerate}[label=\alph*)]
    \item Calculer les premières valeurs.
    \item Conjecturer une expression de \(u_n\).
    \item Démontrer cette expression par récurrence.
\end{enumerate}

\item Écrire un algorithme qui détermine le plus petit entier \(n\) tel que
\[
\frac{n}{2^n}<10^{-4}.
\]

\item On considère une suite définie par récurrence.  
Proposer :

\begin{enumerate}[label=\alph*)]
    \item un algorithme qui calcule \(u_n\) ;
    \item un algorithme qui affiche le premier rang où \(u_n>100\).
\end{enumerate}

\item On veut approcher une solution de l'équation
\[
x^3+x-1=0
\]
sur \([0,1]\).

\begin{enumerate}[label=\alph*)]
    \item Montrer qu'il existe une unique solution sur \([0,1]\).
    \item Décrire une méthode de dichotomie.
    \item Écrire un algorithme associé.
\end{enumerate}

\item On considère une fonction \(f\) strictement croissante sur \([a,b]\).

\begin{enumerate}[label=\alph*)]
    \item Expliquer pourquoi l'équation \(f(x)=0\) admet au plus une solution.
    \item Proposer une méthode algorithmique pour l'approcher si \(f(a)<0<f(b)\).
\end{enumerate}

\end{enumerate}

% =========================================================
\begin{chapitrebox}{Chapitre 12 --- Arithmétique}
\end{chapitrebox}

\begin{enumerate}[label=\textbf{Exercice \arabic*.}, leftmargin=*, itemsep=1em]

\item Déterminer le reste de la division euclidienne de \(2^{10}\) par 7.

\item Montrer que si \(n\) est impair, alors \(n^2\) est impair.

\item Montrer que pour tout entier \(n\),
\[
n(n+1)
\]
est pair.

\item Déterminer les entiers \(n\) tels que
\[
n^2-n
\]
soit divisible par 2.

\item Montrer que pour tout entier \(n\),
\[
n^3-n
\]
est divisible par 6.
\end{enumerate}

\vspace{1em}

\begin{objectifbox}
\textbf{Conseil de travail.} Pour chaque exercice, on attend une rédaction claire, structurée et justifiée.  
Même lorsqu'un calcul semble évident, il faut expliquer le raisonnement.
\end{objectifbox}

\end{document}